\Askhsh{20657}{4}
\begin{ylh}{1}
\item 5.1 Εκθετική συνάρτηση
\item 5.2 Λογάριθμοι
\end{ylh}
Σύμφωνα με τον νόμο ψύξης του Νεύτωνα, η θερμοκρασία $\theta$, σε βαθμούς Κελσίου, ενός αντικειμένου μειώνεται με την πάροδο του χρόνου $t$, σε λεπτά, σύμφωνα με τη συνάρτηση $\theta(t) = T + (\theta_{0} - Τ)e^{kt}$, όπου $k$ μια σταθερά με $k < 0$, $\theta_{0}$ η αρχική θερμοκρασία του αντικειμένου, ενώ $Τ$ είναι η σταθερή θερμοκρασία του περιβάλλοντος μέσα στο οποίο τοποθετείται το αντικείμενο, με $\theta_{0} > Τ$.
Ένα αντικείμενο έχει θερμανθεί στους $100^{ο}\ C$ και στη συνέχεια αφήνεται να κρυώσει σε ένα δωμάτιο με σταθερή θερμοκρασία $30^{ο}\ C$. Γνωρίζουμε ότι 5 λεπτά μετά την τοποθέτησή του αντικειμένου στο δωμάτιο, η θερμοκρασία του αντικειμένου είναι $80^{ο}\ C$.
\begin{alist}
\item Να αποδείξετε ότι $k = - 0,0672$.
\monades{9}
\item Να αποδείξετε ότι $\theta(t) = 30 + 70 \cdot \left( \frac{5}{7} \right)^{\frac{t}{5}}$.
\monades{8}
\item Να βρείτε, με προσέγγιση εκατοστού, τη θερμοκρασία του αντικειμένου μετά από 1 ώρα και 40 λεπτά.
\monades{8}
Δίνεται ότι $\ln\left( \frac{5}{7} \right) = - 0,336$ (προσεγγιστικά) και $\left( \frac{5}{7} \right)^{10} \cong 0,034$.
\end{alist}

